% Programming in Prolog - não esquecer de citá-la

Na experiência do autor do livro (citar livro), Prolog é uma linguagem que
programadores iniciantes acham ser mais compreensiva que programas equivalentes
em linguagens ``convencionais''. Em contra-partida, programadores experientes
nessas mesmas linguagens, apesar de estarem mais bem preparados para os
conceitos abstratos que sondam as LPs, podem achar Prolog uma linguagem díficil
de se acostumar, e precisam de muito convencimento até considerarem a linguagem
como uma ferramenta de programação útil.

Desde seus começos ao redor de 1970, Prolog foi escolhido por muitos
programadores para aplicações de computação simbólica, como:

\begin{itemize}
    % verificar se estes são mesmo os termos em PT-BR
    % ter certeza também que sabe, ao menos em conceito o que significa
    \item Lógica matemática
    \item Resolução de problemas abstratos
    \item Entender linguagem natural
    \item banco de dados relacional
    \item Implementar automação
    \item Resolução de equações simbólicas
    \item Análise de estrutura bioquímica
    \item E várias áreas de inteligência artificial
\end{itemize}

% importante citar também como a linguagem dá nome aos seus elementos sintáticos
% e semânticos, pode-se encontrar isso no livro na seção 1.2 e 1.3 termos
% imporantes: fatos, perguntas, variáveis, conjunções (conjunctions), regras
% listas e recursão

Prolog permite que um computador seja usado como um armazem de fatos e regras,
nos permite a fazer inferências de um fato para outro, encontrar os valores da
variáveis pode nos levar a uma dedução lógica.

\section{Comentários}
\begin{minted}{prolog}
    % esse é um comentário de uma linha só (lembra LaTeX)
    mãe(roberta, joaquina) % roberta é mãe de joaquina

    /*
    rosas são vermelhas
    violetas são azuis
    comentário multilinha
    adoro cuzcuz
    */
\end{minted}

\section{Fatos}

Para expressar a relação: ``João gosta de Maria'', escrevemos em Prolog:

\begin{minted}{prolog}
    gosta(joão, maria).
\end{minted}

A ordem dos fatores é arbitrária, mas deve respeitar alguma ordem determinada
pelo programador, para manter consistência. No exemplo, levemos em consideração
que a pessoa que gosta, seja o primeiro ``elemento'' do que está entre
parênteses. Que não é a mesma coisa que:

\begin{minted}{prolog}
    gosta(maria, joão).
\end{minted}

Para demonstrar com mais clareza ainda como a consistência é importante na
linguagem, observe a seguinte definição de um \emph{fato}:
\begin{minted}{prolog}
    pai(tom, william).
\end{minted}
Podemos assumir que \emph{tom} é pai de \emph{william}, mas também podemos
interpretar que os dois possuem o mesmo pai. A linguagem não restringe a
semântica, que por sua vez, é responsabilidade do programador.

Outros exemplos, com as possíveis interpretações:
\begin{minted}{prolog}
    valioso(diamante).          % Diamante é valioso.
    mulher(jane).               % Jane é uma mulher.
    possui(jane, diamante).     % Jane possui diamantes
    pai(alberto, maria).        % Alberto é pai de Maria
    da(alberto, livro, maria).  % Alberto deu um livro para Maria
\end{minted}

Note que cada fato descrito acima, termina com um ponto final.

Vamos agora dar as terminologias corretas para essa construção; os
\emph{objetos} entre parênteses (de vermelho) são os argumentos, e o nome da
relação que vem antes (de azul), é chamado de \emph{predicado}.

O coletivo de fatos é chamado de \emph{base de dados}.

\section{Perguntas}